% =====================================================================
% 02_standards_spectrum.tex
% =====================================================================
\begingroup
\newcolumntype{L}[1]{>{\hsize=#1\hsize\raggedright\arraybackslash}X}
\newcolumntype{C}[1]{>{\hsize=#1\hsize\centering\arraybackslash}X}

\section{Tiêu chuẩn IEEE 802.11 và phổ tần}
\label{sec:standards}

\subsection{Các thế hệ Wi-Fi}

Bảng~\ref{tab:generations} tóm tắt các phiên bản sửa đổi chính của lớp PHY.
Tên thế hệ dành cho 802.11b/a/g là cách gọi không chính thức, được đặt ra về sau.

\begin{table}[htbp]
  \centering\footnotesize
  \caption{Các thế hệ Wi-Fi và phiên bản sửa đổi PHY.}
  \label{tab:generations}
  \begin{tabularx}{\linewidth}{@{} L{1.1} C{0.4} C{0.6} C{0.5} C{0.4} C{0.8} L{2.2} @{}}
    \toprule
    \textbf{Thế hệ / phiên bản sửa đổi} & \textbf{Năm} & \textbf{Băng tần (GHz)} & \textbf{Độ rộng tối đa (MHz)} & \textbf{QAM tối đa} & \textbf{Tốc độ PHY đỉnh} & \textbf{Công nghệ chính} \\
    \midrule
    -- / 802.11                     & 1997          & 2.4                  & 22                       & --               & 2~Mb/s                 & \acs{fhss}, \acs{dsss}, hồng ngoại; \acs{csmaca} \\
    -- / 802.11b                    & 1999          & 2.4                  & 22                       & --               & 11~Mb/s                & HR-DSSS với \acs{cck} \\
    -- / 802.11a                    & 1999          & 5                    & 20                       & 64               & 54~Mb/s                & \acs{ofdm} (52 sóng mang con) \\
    -- / 802.11g                    & 2003          & 2.4                  & 20                       & 64               & 54~Mb/s                & \acs{ofdm} ở 2.4~GHz, tương thích với 11b \\
    Wi-Fi~4 / 802.11n               & 2009          & 2.4, 5               & 40                       & 64               & 600~Mb/s               & \acs{mimo} (4 luồng), ghép kênh, \acs{ampdu} \\
    Wi-Fi~5 / 802.11ac              & 2013          & 5                    & 160                      & 256              & 6.93~Gb/s              & 8 luồng, \acs{mumimo} đường xuống, tạo búp sóng (beamforming) \\
    Wi-Fi~6/6E / 802.11ax           & 2021          & 2.4, 5, 6            & 160                      & 1024             & 9.6~Gb/s               & \acs{ofdma}, \acs{mumimo} lên/xuống, tô màu BSS, \acs{twt} \\
    Wi-Fi~7 / 802.11be              & 2024          & 2.4, 5, 6            & 320                      & 4096             & khoảng 46~Gb/s         & \acs{mlo}, nhiều \acs{ru}, loại bỏ một phần kênh (puncturing) \\
    \bottomrule
  \end{tabularx}
\end{table}

\subsection{Các phiên bản sửa đổi chức năng}

Bên cạnh các phiên bản sửa đổi PHY, nhiều phiên bản khác bổ sung chức năng ở
lớp MAC. Những phiên bản tiêu biểu được liệt kê trong
Bảng~\ref{tab:amendments}.

\begin{table}[htbp]
  \centering\footnotesize
  \caption{Các phiên bản sửa đổi chức năng của 802.11.}
  \label{tab:amendments}
  \begin{tabularx}{\linewidth}{@{} L{0.5} L{0.9} L{1.6} @{}}
    \toprule
    \textbf{Phiên bản sửa đổi} & \textbf{Chủ đề} & \textbf{Mục đích} \\
    \midrule
    802.11e            & \acs{qos}                  & Truy cập kênh theo mức ưu tiên (\acs{edca}), nền tảng của \acs{wmm} \\
    802.11h            & Quản lý phổ tần            & \acs{dfs} và \acs{tpc} để cùng tồn tại với radar ở 5~GHz \\
    802.11i            & Bảo mật                    & \acs{ccmp}/\acs{aes}, bắt tay 4 bước (WPA2) \\
    802.11k            & Đo lường tài nguyên vô tuyến & Báo cáo mạng lân cận và phép đo beacon hỗ trợ roaming \\
    802.11r            & Chuyển tiếp BSS nhanh       & Giảm độ trễ chuyển vùng nhờ dẫn xuất khóa trước \\
    802.11v            & \acs{btm}                  & Mạng có thể đề xuất \acs{ap} phù hợp hơn cho thiết bị khách \\
    802.11w            & \acs{pmf}                  & Xác thực khung hủy xác thực và ngắt liên kết \\
    802.11s            & Mạng mesh                  & Chuyển tiếp nhiều chặng giữa các \acp{ap} mesh \\
    802.11u            & Kết nối liên mạng           & Nền tảng cho thỏa thuận roaming Passpoint/Hotspot~2.0 \\
    802.11ai           & \acs{fils}                 & Xác thực ban đầu và thiết lập liên kết nhanh \\
    802.11ad / ay      & 60~GHz (DMG / EDMG)        & Liên kết đa gigabit trong phạm vi ngắn \\
    802.11ah           & Dưới 1~GHz (Wi-Fi HaLow)   & \ac{iot} tầm xa, công suất thấp \\
    802.11az           & Định vị                    & Cải thiện đo khoảng cách và độ chính xác vị trí \\
    802.11ba           & Vô tuyến đánh thức         & Đánh thức thiết bị \ac{iot} với công suất cực thấp \\
    802.11bf           & Cảm biến WLAN              & Dùng tín hiệu Wi-Fi để phát hiện sự hiện diện và chuyển động \\
    \bottomrule
  \end{tabularx}
\end{table}

\subsection{Băng tần và kênh}

\subsubsection{Băng tần 2,4~GHz}

Dải ISM 2,4~GHz trải từ khoảng 2400 đến 2483,5~MHz, có tầm phủ và khả năng
xuyên tường khá tốt nhưng băng thông hẹp và thường đông đúc do cùng được sử
dụng bởi Bluetooth, Zigbee, lò vi sóng và điện thoại không dây. Các kênh cách
nhau 5~MHz, trong khi tín hiệu rộng 20/22~MHz chồng lấn lên khoảng bốn số kênh
kề nhau. Vì vậy, kênh 1, 6 và 11 thường được chọn để hạn chế chồng lấn \cite{cisco_rf_guide}. Tại
Việt Nam, băng tần này sử dụng các kênh 1--13; thiết bị phải tuân thủ quy chuẩn
kỹ thuật của \ac{vnmic}, chẳng hạn \acs{qcvn} 54:2020/BTTTT, về công suất phát
cho phép.

\begin{figure}[htbp]
  \centering
  \begin{tikzpicture}
    \pgfmathsetmacro{\xscale}{0.83}
    \pgfmathsetmacro{\wband}{1.43}
    \pgfmathsetmacro{\wctrl}{0.6}
    \pgfmathsetmacro{\hpeak}{1.6}

    \foreach \n in
    {1,...,13}{
      \pgfmathsetmacro{\xc}{(\n - 1) * \xscale} \draw[gray!70] (\xc-\wband,0)
      .. controls (\xc-\wctrl,\hpeak) and (\xc+\wctrl,\hpeak) ..
      (\xc+\wband,0); \node[font=\scriptsize, anchor=north] at (\xc,-0.1)
      {\n};
    }
    \foreach \n/\col in
    {1/red, 6/blue, 11/green!50!black}{
      \pgfmathsetmacro{\xc}{(\n - 1) * \xscale} \draw[\col,very thick]
      (\xc-\wband,0) .. controls (\xc-\wctrl,\hpeak) and (\xc+\wctrl,\hpeak)
      .. (\xc+\wband,0);
    }
    \pgfmathsetmacro{\xend}{12 * \xscale + 1.2} \draw[->] (-1.5,0)
    -- (\xend,0);
    \node[font=\small, anchor=north] at ({6 * \xscale},-0.6)
    {Số kênh (tần số trung tâm $= 2407 + 5n$~MHz)};
  \end{tikzpicture}
  \caption{Quy hoạch kênh 2.4~GHz: các kênh 1, 6 và 11 (nét đậm) không chồng lấn.}
  \label{fig:ch24}
\end{figure}

\subsubsection{Băng tần 5~GHz}

Băng tần 5~GHz được chia thành các dải nhỏ với điều kiện sử dụng trong nhà,
ngoài trời và giới hạn công suất khác nhau. Một số kênh yêu cầu \ac{dfs}:
\ac{ap} phải dò tín hiệu radar và chuyển khỏi kênh nếu phát hiện radar. Tại
Việt Nam, việc triển khai WLAN ở băng 5~GHz chịu sự quản lý chặt chẽ của
\acs{vnmic}; thiết bị phải đáp ứng các quy chuẩn \acs{qcvn} tương ứng để được
chứng nhận hợp quy trước khi thông quan hoặc đưa vào sử dụng. Tần số trung tâm
của kênh 5~GHz số $n$ được tính như sau:
\begin{equation}
  f_c = 5000 + 5n \ \text{MHz}.
\end{equation}
Các kênh có thể được ghép thành kênh rộng 40, 80 hoặc 160~MHz. Tuy nhiên, kênh
càng rộng thì số kênh độc lập càng ít và nguy cơ nhiễu đồng kênh càng cao.

\subsubsection{Băng tần 6~GHz}

Wi-Fi~6E và Wi-Fi~7 có thể hoạt động trên băng tần 6~GHz. Theo Thông tư
01/2025/TT-BKHCN của Việt Nam (có hiệu lực từ ngày 15~tháng 5~năm 2025), dải
5925--6425~MHz được phân bổ cho Wi-Fi. Mức sử dụng trong nhà được cho phép đến
200~mW \ac{eirp} (mật độ công suất phổ \ac{psd} tối đa 10~dBm/MHz); triển khai
ngoài trời bị giới hạn ở 25~mW \ac{eirp} (\ac{psd} tối đa 1~dBm/MHz). Không
được lắp đặt thiết bị trên các hệ thống máy bay không người lái (\acp{uas}).
Ở nơi được phép sử dụng băng tần 6~GHz, bắt buộc dùng WPA3 (hoặc \ac{owe}).

\subsubsection{Các băng tần khác}

Băng tần 60~GHz (802.11ad/ay) cung cấp tốc độ rất cao trong phạm vi vài mét,
thường giới hạn trong một phòng. Băng tần dưới 1~GHz (802.11ah, ``HaLow'')
đánh đổi tốc độ để tăng tầm phủ và khả năng xuyên tường, phù hợp với cảm biến
\ac{iot}.

\subsection{Quy định quản lý tần số}

Tại Việt Nam, \ac{vnmic} ban hành quy hoạch phân bổ tần số và quy chuẩn kỹ
thuật quốc gia (\acs{qcvn}), quy định các kênh WLAN được phép sử dụng, công
suất phát, điều kiện triển khai trong nhà/ngoài trời và yêu cầu \ac{dfs}. Thiết
bị phải được thử nghiệm, chứng nhận hợp quy và nộp bản công bố hợp quy (\ac{doc})
theo quy định hiện hành.

\begin{tomtat}
\item Wi-Fi phát triển từ 802.11 (1997, 2~Mb/s) đến Wi-Fi~7 (802.11be: kênh 320~MHz, 4096-QAM, \acs{mlo}, tốc độ \acs{phy} lý thuyết tối đa khoảng 46~Gb/s). Tốc độ đỉnh chỉ là lý thuyết.
\item Các bản bổ sung chức năng cho \acs{phy}: \acs{qos} (e), \acs{dfs}/\acs{tpc} (h), bảo mật (i, w), chuyển vùng (k, r, v), mesh (s), cùng các bản cho \acs{iot}, định vị và cảm biến.
\item Băng 2,4~GHz: tầm xa, xuyên tường tốt nhưng hẹp và đông đúc; kênh 1, 6, 11 không chồng lấn; Việt Nam dùng kênh 1--13.
\item Băng 5~GHz: nhiều kênh hơn, có thể ghép 40/80/160~MHz, một số kênh yêu cầu \acs{dfs} để tránh radar.
\item Băng 6~GHz: theo Thông tư 01/2025/TT-BKHCN, dải 5925--6425~MHz dành cho Wi-Fi; trong nhà tối đa 200~mW \acs{eirp}, ngoài trời 25~mW, không được sử dụng trên \acs{uas}; bắt buộc WPA3 hoặc \acs{owe}.
\item Thiết bị tại Việt Nam phải tuân thủ các quy chuẩn \acs{qcvn} do \acs{vnmic} quy định, được kiểm tra hợp quy và nộp bản \acs{doc}.
\end{tomtat}
\endgroup
